\chapter*{Abstract}

\noindent Much of the clinical information on patients with liver disease is recorded as free text in electronic health records, which makes its manual extraction slow and its systematic analysis difficult. Large language models can automate this extraction, but, when combined with traditional similarity-based retrieval, they struggle with evidence scattered across the text. This work evaluates GraphRAG, an approach that combines these models with knowledge graphs, for identifying eight liver diseases in real Brazilian Portuguese electronic health records from the University Hospital of the Federal University of Alagoas. On a sample of 537 manually annotated records, the \textit{basic}, \textit{local} and \textit{drift} search modes were compared at two community levels of the graph. The \textit{drift} mode performed best (macro F1 of 0.798), owing to a much higher recall at the cost of lower precision. The gain was largest for diseases whose evidence is scattered across the text, such as cirrhosis and portal hypertension, and small or absent for those with a direct textual marker; the community level mattered less than the search mode. We conclude that the choice of search mode should take into account the type of textual evidence available for each disease.

\vspace{5mm}
\noindent\textbf{\textit{Keywords}:~}GraphRAG; natural language processing; electronic health records; hepatology; multilabel classification.
